\documentclass[11pt]{article}

\usepackage[T1]{fontenc}
\usepackage[utf8]{inputenc}
\usepackage{amsmath,amssymb}
\usepackage[margin=1in]{geometry}
\usepackage{xcolor}
\usepackage{graphicx}
\usepackage{booktabs}
\usepackage{longtable}
\usepackage[round,authoryear]{natbib}
\usepackage[hidelinks]{hyperref}
\graphicspath{{paper/old_paper_ml_structure/figures/}}
\AddToHook{env/figure/begin}{\centering}
\AddToHook{env/table/begin}{\centering}
\setlength{\LTleft}{\fill}
\setlength{\LTright}{\fill}

\definecolor{boxrule}{HTML}{2C6EAB}
\newcommand{\fit}{\phi}
\newcommand{\data}{\mathcal{D}}
\newcommand{\resolved}{\rho_{\alpha}}
\newcommand{\rank}{\mathrm{rank}}
\newcommand{\note}[1]{\textcolor{boxrule}{\small[#1]}}

\title{How well do deep mutational scans resolve top-ranked variant batches?\\
\large Ranking uncertainty in protein engineering and machine-learning benchmarks}
\author{[TODO: authors and affiliations]}
\date{Working draft for bioRxiv \textbar\ September 2026}

\begin{document}
\maketitle

\noindent\note{Working draft: Numerical results have been checked against the
saved analysis tables. Outstanding references, fitting details, and publication
metadata are marked TODO.}

\begin{abstract}
Protein engineers use measurements of protein fitness, performance in a
specified experimental test, to choose batches of variants for further study.
Machine-learning researchers use the same measurements to train and evaluate
predictors of fitness from sequence. Deep mutational scans provide these
measurements for many variants, and sorting the reported scores gives a
convenient ranking. But the selected batch depends on both the calculation
of fitness from sequencing counts and uncertainty in those counts. In the
present work, we quantify these two sources of variation in the leading
ranks of three protein systems. We fit models to the counts, generate
sampled fitness scores, and measure rank intervals, how often each candidate
is selected, and overlap between batches. Changing the fitness calculation
changes batch membership in the examples examined, and noise within a fixed
analysis can replace a substantial fraction of a batch. The extent of this
turnover differs across landscapes and selection budgets, so an exact list
of leading variants can conceal substantial uncertainty about which
candidates to select.
\end{abstract}

\section{Introduction}
\label{sec:introduction}

Protein engineers want to identify variants with high protein fitness:
good performance in a specified experimental test \citep{yang2019}.
For example, an engineer may select a hundred variants for validation,
while a machine-learning researcher asks whether a predictor identifies
the same candidates. Deep mutational scans measure many variants, sequences
that differ from a reference protein, in parallel \citep{fowler2014}.
Their fitness scores guide batch selection and serve as labels for learning
to predict fitness from sequence. Benchmarks such as FLIP, ProteinGym,
and FLIP2 evaluate predictions against these measurements
\citep{dallago2021,notin2023,didi2026}.

A released \emph{fitness landscape} is usually a table of sequences and
scores. Sorting those scores gives the \emph{reported ordering}, and taking
the highest-scoring variants gives a candidate batch. This is a natural
way to use the experiment, but it leaves two questions unanswered. First,
would another calculation of fitness select the same variants? Fitness
depends on the test, or \emph{assay}, and on choices such as normalization
and how repeated measurements are combined. Second, how certain is the
selection under a fixed calculation? In these experiments, sequencing
counts how often each variant is observed before and after selection or
sorting. Those finite counts leave uncertainty in the fitness estimates,
which an exact sorted list does not display.

Methods for estimating fitness already account for measurement uncertainty.
Enrich2 and DiMSum estimate errors from counts and repeated measurements
\citep{rubin2017,faure2020}; model-based enrichment estimation shares
information between related sequences \citep{busia2023}. Experimental-design
analyses also study how read depth and sampling choices affect precision
\citep{matuszewski2016}. But an error estimate for each score does not
tell an analyst how many members of a selected batch could change. That
depends on the competing scores and how their uncertainties vary together.
We use \emph{ranking resolution} to describe uncertainty in the ranks and
batch membership of leading variants. Our question is how much the score
uncertainty changes these ranking and selection conclusions.

In the present work, we study scans of three protein systems: protein
binding through PDZ interactions, tryptophan
production by TrpB, and DNA cleavage by NucB
\citep{zarin2024,johnston2024,thomas2025}. We fit statistical models directly
to the sequencing counts and sample fitness scores or alternative sequencing
outcomes under those models. We rank each sampled set of scores and select
its leading batch. We compare reported and fitted orderings
(Section~\ref{sec:reported}), measure rank uncertainty and batch turnover
within each analysis (Section~\ref{sec:noise}), and examine their relationships
with count support and competing scores (Section~\ref{sec:mechanisms}).
The reported top 100 and top 1,000 and batch sizes from 10 to 1,000 define
the scope of these comparisons.

Changing the fitness calculation changes the selected batch in the examples
we examine. Within a fixed analysis, noise can also replace a substantial
fraction of its members. The amount of turnover differs across landscapes
and does not always decrease as batch size grows. Lower input counts are
associated with wider score intervals among the leading candidates in every
analysis.
These results quantify uncertainty in the experimental lists used for
engineering and prediction benchmarks, before introducing a fitness predictor.
The landscapes studied here contribute to FLIP2; evaluating the effect on
its particular predictor comparisons remains future work.
\note{TODO: Add versioned code and data links once their public availability
is confirmed.}

\section{From sequencing counts to ranking decisions}
\label{sec:background}

\paragraph{Fitness depends on the experiment.}
A protein is a chain of amino acids, and its sequence helps determine what
it binds or which chemical reactions it helps carry out. Changing the sequence
can affect these properties differently. A test of binding to another protein
and a test of enzyme activity therefore ask different questions about the
same variant. There is no single assay-independent fitness that both must
recover. In this paper, fitness is the score defined by the experimental
test and the calculation applied to its measurements
\citep{fowler2014,johnston2024}. A high score means good performance according
to that test and calculation, not universal usefulness.

\paragraph{From variants to sequencing counts.}
The experiments begin with a \emph{library}, a collection of DNA sequences
encoding different protein variants. Researchers test these variants together
and connect each protein's performance to the DNA that identifies it.
First, sequencing samples the library before selection, giving an input
count for each variant. The experiment then favors variants according to
the property being tested. Finally, sequencing samples the selected population,
giving output counts. Comparing input and output accounts for the fact that
variants need not have been equally common at the start.

The physical selection differs between assays. In a growth selection, cells
carrying different variants reproduce at rates affected by the tested protein
property. In a fluorescence screen, a signal reports activity and a sorting
device collects cells or droplets above a chosen signal threshold, called
a \emph{gate}. A study can sequence several timepoints or several sorted
populations, rather than just one output. The next section describes the
particular assays used here.

A sequencing \emph{read} is an observation of a DNA fragment assigned to
a variant. The variant's count is the number of reads assigned to it, and
the total assigned reads in a sequenced population are its sequencing depth.
A read identifies a variant; it does not directly measure that protein's
binding strength or reaction rate. The data used to calculate fitness are
therefore counts, not observations of an exact score or rank.

\paragraph{From counts to a reported score.}
A common starting point is \emph{enrichment}: the fraction of output reads
assigned to a variant divided by its fraction of input reads. Fractions
account for differences in the total number of reads collected before and
after selection. The calculation may then compare the variant with a
reference sequence, combine repeated measurements, or summarize changes
across timepoints. The unmutated reference sequence is often called wild
type; other known sequences can serve as experimental controls. These choices
define the score and can affect the resulting ranking.

\paragraph{Why counts leave uncertainty.}
The ratio alone hides how much data supports it. Consider a hypothetical
comparison with equal input and output sequencing depths: one variant has
10 output reads for 1 input read, and another has 10,000 output reads for
1,000 input reads. Both ratios are ten. But one additional input read
halves the first ratio and barely changes the second. The same estimated
enrichment can therefore come with very different uncertainty. A large
experiment can contain both well-observed and sparsely observed variants.

Zero counts require additional choices because division by zero is undefined
and an observed zero need not mean a variant was absent from the population.
For example, adding a small count before taking a ratio or excluding a
poorly observed variant changes the calculation. When only the final score
is used downstream, neither these choices nor the amount of supporting
data is apparent. Enrich2 and DiMSum retain uncertainty when estimating
fitness from counts \citep{rubin2017,faure2020}; we examine how that
uncertainty affects ranks and selected batches.

\paragraph{Repeated measurements.}
Repeating the biological experiment with another population gives a
biological replicate. Sequencing the same population again instead measures
variation in the sequencing step. Later timepoints and different sorting
gates provide other observations, but do not repeat the same experiment.
These distinctions matter when combining counts: disagreement between
observations can reflect both limited sampling and differences in the
populations or conditions. The models retain the observation structure
rather than treating every count vector as a biological replicate.

\paragraph{From scores to ranks and batches.}
A score describes a variant on the experiment's scale; a rank locates it
relative to the other candidates. A top-100 batch contains the hundred
highest-scoring candidates, regardless of their order within that batch.
Two candidates can exchange positions many times while both remain selected.
A smaller movement near the selection cutoff can instead replace a member.
We therefore examine batch membership separately from internal ordering.

It also matters how scores vary together. If every sampled score increases
by the same amount, none of the ranks changes. If scores are close and vary
relative to each other, even small changes can alter their order. Separate
error bars do not describe this difference. We keep the sampled scores for
all variants together and rank each complete set, so the analysis retains
their shared variation.

\paragraph{The decisions we study.}
We compare the reported top 100 and top 1,000 with the fitted analyses,
and examine batch sizes from 10 to 1,000. The top 100 is the main worked
example, not a universal experimental budget. Detailed displays include
members of the proposed batch and outside candidates that could plausibly
replace them; Section~\ref{sec:methods} defines these groups precisely.
Restricting the display does not remove competitors from the ranking:
each displayed variant is still ranked against all available competitors
in that sample.

\section{Assay data and uncertainty analysis}
\label{sec:methods}

\paragraph{Data and scope.}
We analyze five PDZ landscapes, nine TrpB libraries with three varied sites,
one TrpB library with four varied sites, and four NucB engineering generations.
For each landscape, we have variant sequences, reported fitness scores, and
sequencing counts indexed by replicate, timepoint, or sorting gate. We fit
the counts and obtain sampled scores for all candidates together. Ranking
these samples gives a distribution of ranks and selected batches for each
analysis. The assays differ in what fitness measures and how selection
produces the output counts.

\paragraph{PDZ: binding measured through growth.}
The PDZ data measure binding between two proteins \citep{zarin2024}.
PDZ3 is a domain, a structural unit
within the human protein PSD95, that binds a short amino-acid segment of
the protein CRIPT. The source study changes amino acids in the binding
partners and measures how those changes affect their interaction. A candidate
in a joint library therefore specifies a pair of interacting sequences,
not necessarily a mutation in PDZ3 alone.

The experiment couples binding to yeast growth. In a
protein complementation assay, the binding partners bring complementary
reporter fragments together; productive interaction supports growth under
selection with the drug methotrexate. Cells carrying different variants compete
in the same culture. Sequencing before and after this competition measures
which variants increased in relative abundance. The reported DiMSum score
summarizes that change across three biological replicates. It is a
binding-dependent growth score, not a direct measurement of a binding
dissociation constant. Protein abundance, folding, and binding can all
affect this readout. We retain the source labels cis,
trans-1, trans-2, CRIPT-C, and CRIPT-N for the five prepared landscapes;
Appendix~\ref{app:experimental-records} records the available library details.

\paragraph{TrpB: catalysis measured through growth.}
TrpB is the enzyme subunit that makes the amino acid tryptophan from indole
and serine. \citet{johnston2024} express TrpB variants in bacteria whose
native tryptophan-synthesis genes have been deleted. With indole supplied,
the introduced enzyme's production of tryptophan supports bacterial growth.
Cells carrying active variants can increase in the pooled culture, and
sequencing tracks their abundance at the start and at later timepoints.

The three-site libraries A--I vary different triples of amino-acid positions
in an engineered TrpB background; they are different candidate collections,
not replicate measurements of one collection. We also analyze a four-site
library. The reported fitness estimates summarize growth inferred from
abundance changes, with normalization using nonfunctional stop-codon variants
and the highest estimated growth. Our models instead fit effective growth
parameters to the time-indexed counts. Neither score is a purified enzyme's
catalytic rate constant: growth also depends on the amount of functional
enzyme in the cell and the assay environment. The prepared three-site
libraries have one shared input count vector and library-specific output
replicate structures; the four-site landscape has two input observations.

\paragraph{NucB: catalysis measured through fluorescence sorting.}
NucB is a nuclease, an enzyme that cleaves DNA. \citet{thomas2025} engineer
it through four successive generations of designed and screened variants.
Here fluorescence reports nuclease activity, and sorting collects populations
whose signals exceed chosen thresholds. Unlike PDZ and TrpB, the selection
readout is passage through an activity gate, not growth during a competition.
Sequencing the unsorted input and each selected population identifies how
strongly each variant is enriched by that gate. A variant can have many
output reads because it was common initially, because it was more likely
to pass the gate, or both; the input counts help distinguish these cases.

The source analysis uses enrichment relative to control sequences to assess
activity. We retain enrichment at the highest selected gate as the reported
score and fit the counts from all available gates. The fitted score describes
the location of a variant's fluorescence distribution
(Appendix~\ref{app:nuc-model}). Generations G1--G4 contain different candidate
collections and screening protocols, so we analyze them separately.
A \emph{fiducial} is a known control sequence included in the experiment.
The appended premature-stop fiducials encode truncated proteins; we retain
their counts where required for normalization but exclude them from candidate
rankings. The fitted count model can therefore include controls that are
excluded from ranking.

\paragraph{Count models.}
Let $\data$ denote the observed counts and their replicate, timepoint,
or gate identities. Write $\theta$ for the model parameters and
$p(\data\mid\theta)$ for the count likelihood. Each sequenced population
has a vector of variant probabilities. Conditional on those probabilities
and the recorded total number of reads, its count vector follows a
multinomial distribution. This retains the dependence between variant counts
that share a sequencing total. Assay-specific models determine how input
probabilities and protein performance change the output probabilities.

\paragraph{Growth models.}
For PDZ and TrpB, exponential or sigmoid growth functions map input
abundances and variant-specific rates to expected output abundances. The
Exponential and Lomax specifications assign different priors to a positive
central growth rate. The Gamma specification assigns a Gamma prior to an
inverse-rate parameter, which is the rate of the exponential distribution
for replicate growth rates. We name each specification rate-family--growth-function,
such as Exponential--Exponential. Appendices~\ref{app:pdz-model}
and~\ref{app:trpb-model} give the equations and hyperparameters.

PDZ has three replicate-level rate vectors linked by a shared variant-level
parameter. TrpB libraries A--C use one rate per variant across their observed
times; libraries D--I and the four-site model allow replicate-specific rates.
TrpB observation times are divided by the last modeled time, so these growth
parameters describe change over the modeled assay duration rather than per hour.

\paragraph{Fluorescence models.}
For NucB, the model assigns each variant a fluorescence location and a
distribution of observed signals around it. A gate likelihood combines the
input composition with the probability that the signal exceeds the gate's
threshold, then normalizes across variants before evaluating a multinomial
count likelihood. The Exponential and Lomax specifications differ in their
prior on the positive fluorescence-scale location. Generation-specific input
populations and gate thresholds determine which observations share parameters.
Appendix~\ref{app:nuc-model} defines those probabilities, the normal and
lognormal signal options, and all threshold values.

\paragraph{Fitness scores.}
Let $N$ be the number of eligible ranking candidates and
$i\in\{1,\ldots,N\}$ index them. We retain the prepared candidate set for
each analysis, including the source-aligned PDZ and TrpB candidates and
excluding the appended NucB controls. Candidate counts and exclusions are
recorded in Appendix~\ref{app:reproducibility}.

The score $h_i(\theta)$ ranks variant $i$ from a fitted parameter vector.
For TrpB models without replicate effects, $h_i=\log b_i$, where $b_i>0$ is the variant's
growth rate. For Exponential or Lomax growth hierarchies, $h_i=\log m_i$,
where $m_i$ is the replicate-rate center. It is a conditional mean for normal
replicate variation and a conditional median for lognormal variation.
For the Gamma inverse-rate hierarchy, $h_i=-\log\kappa_i$, where
$1/\kappa_i$ is the conditional mean replicate rate. For NucB, $h_i=L_i$,
where $L_i$ is the fitted log-fluorescence location. Higher scores are
preferred in all three assays, but the units differ between protein systems.

We also calculate scores directly from counts, denoted $g_i(\data)$. The
PDZ calculation combines replicate log enrichments relative to wild type
using inverse-variance weights. It approximates, rather than reruns, the
published DiMSum analysis. The TrpB calculation uses the deposited stop-variant
and maximum-score normalization before averaging timepoints and replicates.
The NucB calculation uses the highest selected gate's enrichment. The TrpB
and NucB calculations reproduce the source scores from observed counts;
Appendix~\ref{app:scores} gives formulas, tolerances, and filtering rules.

\paragraph{Fitting.}
We combine the count likelihood with a prior $p(\theta)$ and draw from
$p(\theta\mid\data)\propto p(\data\mid\theta)p(\theta)$ using Hamiltonian
Monte Carlo (HMC). We also obtain a single parameter vector $\hat\theta$
by gradient-based optimization of the joint density, including likelihood
and prior terms. We call this an \emph{optimized fit}.
Appendix~\ref{app:inference-details} gives the objective and per-fit settings.

\paragraph{Sampling fitness.}
One \emph{run} combines a landscape, model specification, and resampling
scheme. Let $S$ be its number of retained samples, indexed by
$s\in\{1,\ldots,S\}$, and
$\boldsymbol{\fit}^{(s)}=(\fit_1^{(s)},\ldots,\fit_N^{(s)})$ its joint
sampled score vector. We use three recipes. Parameter resampling
draws $\theta^{(s)}\sim p(\theta\mid\data)$ and sets
$\fit_i^{(s)}:=h_i(\theta^{(s)})$. Optimized-fit count resampling draws
$\data^{(s)}\sim p(\cdot\mid\hat\theta)$ and sets
$\fit_i^{(s)}:=g_i(\data^{(s)})$. Posterior count resampling draws
$\theta^{(s)}\sim p(\theta\mid\data)$, then
$\data^{(s)}\sim p(\cdot\mid\theta^{(s)})$, and again sets
$\fit_i^{(s)}:=g_i(\data^{(s)})$. The count-based schemes recompute fitness
without refitting the model. Each score vector retains the dependence
between candidates within one draw.

Parameter draws describe uncertainty in the fitted score. Count draws
describe alternative sequencing outcomes at the recorded depths under the
fitted model. They use the count-derived score and do not reproduce every
source of variation in a new physical experiment. We apply the same ranking
summaries to all three recipes, keeping them separate because they differ
in both score calculation and uncertainty. Table~\ref{tab:inventory} lists
the available analyses; the four-site TrpB landscape has optimized-fit
count samples only.

\begin{table}[t]
\centering
\caption{Analyses used for the ranking summaries. A run is one landscape,
model specification, and resampling scheme. Retained samples are draws per
run, not effective sample sizes. Models fitted to one landscape share data.}
\label{tab:inventory}
\begin{tabular}{lrrr}
\toprule
Scheme & Landscapes & Runs & Retained samples \\
\midrule
Parameter & 18 & 49 & 16,000--32,000 \\
Optimized-fit counts & 19 & 60 & 967--1,000 \\
Posterior counts & 18 & 49 & 15,839--32,000 \\
\bottomrule
\end{tabular}
\end{table}

\paragraph{Ranking summaries.}
Let $R_i^{(s)}$ be the descending rank of $\fit_i^{(s)}$ among
ranking-eligible candidates with finite scores in sample $s$. Rank one is
best, and exact score ties receive their average rank. When the focal
candidate's score is undefined, its rank is also undefined. We write
$Q_a(\fit_i)$ and $Q_a(R_i)$ for empirical quantiles, at probability $a$, of
the finite samples $\{\fit_i^{(s)}\}$ and $\{R_i^{(s)}\}$, respectively.
Thus unsuperscripted scores and ranks inside $Q_a$ refer to sample collections,
with quantiles calculated separately for each candidate.

We define the \emph{median-score reference} by taking
$\widetilde\fit_i:=Q_{0.50}(\fit_i)$ for each candidate and ranking these
medians to obtain $\widetilde r_i$. Ranking median scores need not give
the median sampled ranks. Below, the \emph{reference top 100} means the
hundred candidates with the highest median scores. We summarize rank uncertainty using
$[Q_{0.05}(R_i),Q_{0.95}(R_i)]$, its width in rank positions, and the two
endpoints minus $\widetilde r_i$. These signed deviations describe movement
relative to the median-score reference. Its rank need not lie inside the
sampled-rank interval. The intervals describe the retained sample distributions;
they are not a claim of frequentist coverage.

\paragraph{Candidate selection.}
For a selection budget $k$, let $T_k^{(s)}$ be the set of the $k$ highest-scoring
eligible variants with finite scores in sample $s$. We estimate inclusion
probabilities by $\hat q_i(k):=S^{-1}\sum_{s=1}^S\mathbf{1}\{i\in T_k^{(s)}\}$,
where the indicator equals one if the candidate is selected and zero otherwise.
An undefined score counts as no selection; the denominator includes every
retained sample. Let
$T_k^\star$ and $T_k^{\mathrm{rep}}$ denote the sets induced by the median-score
reference and reported orderings, respectively. We report three shared fractions:
\begin{align}
\widehat C_k
&:=\frac{1}{S(S-1)}\sum_{s\ne s'}
\frac{|T_k^{(s)}\cap T_k^{(s')}|}{k},\\
\widehat C_k^\star
&:=\frac{1}{S}\sum_s\frac{|T_k^{(s)}\cap T_k^\star|}{k},\qquad
\widehat O_k^{\mathrm{rep}}
:=\frac{1}{S}\sum_s\frac{|T_k^{(s)}\cap T_k^{\mathrm{rep}}|}{k}.
\end{align}
The first measures agreement between two resampled selections; the second
compares samples with their median-score reference; the third compares samples
with the reported batch. We examine batch budgets
$k\in\{10,20,50,100,200,500,1000\}$. At these budgets, a shared fraction
is not the frequency of exact set agreement: an overlap of 0.8 for a
top-100 batch means 80 shared members on average, not an 80\% chance that
the complete batches match.
We retain fixed-size sets using the deterministic tie convention of the
stored selections, rather than thresholding the average sampled ranks.
Appendix~\ref{app:computation} describes finite-score rules and the exact
overlap calculation. Table~\ref{tab:reference-comparisons} reports all three
quantities for the main examples, alongside the intersection of the two
fixed reference sets.

\paragraph{Analysis cohorts.}
Reported-order comparisons use the reported top 100 and top 1,000. Detailed
rank-interval summaries use the median-score reference top 100. Our
\emph{decision cohort} contains that reference top 100 and every other
candidate with $\hat q_i(100)\geq0.01$. We call the latter candidates
\emph{challengers}; they can lie outside the reference top 1,000.
A displayed candidate is always ranked
against all finite eligible competitors in its sample, not just against the
other displayed candidates.

\paragraph{Comparisons and aggregation.}
As a reference for set overlap, two independent uniformly
selected size-$k$ subsets of $N$ candidates share fraction $k/N$ in expectation.
We report individual runs and summarize model specifications within landscapes
by their median, so an experiment with more fitted models does not count
as more experimental evidence. Comparisons across resampling schemes use
the same named landscape and model specification where available. Candidate
sets can differ slightly because the score calculations have different
exclusion rules (Appendix~\ref{app:computation}).

\section{Does the fitness construction change the ordering?}
\label{sec:reported}

The reported and fitted scores select different top-100 batches. Under
parameter resampling, the median-score reference shares 52 candidates
(52\%) with the reported top 100 for PDZ trans-1, 33 (33\%) for TrpB library
B, and 68 (68\%) for NucB generation 3 under the example specifications in
Table~\ref{tab:reference-comparisons}. These intersections compare two fixed
lists. They measure disagreement between score calculations before we
consider how the selected batch varies across samples.

\input{paper/old_paper_ml_structure/figures/reference_comparisons.tex}

The score calculations need not estimate the same assay summary. For NucB,
the reported comparator uses the highest selected gate, while the fitted
parameter score uses counts from all available gates. For PDZ, our
count-based score approximates the original DiMSum calculation. Disagreement
therefore does not establish that the reported scores are wrong.

We next compare reported ranks with the sampled rank distributions. These
diagnostics depend on both the fitted scores and their uncertainty. They
are distinct from the fixed-list intersections above.

\paragraph{Reported ranks within sampled distributions.}
For each candidate, we calculate the fraction of finite sampled ranks that
are better than its reported rank, giving ties half weight. A fraction
near zero means that the reported rank is better than most sampled ranks;
a fraction near one means it is worse. Figure~\ref{fig:reported}a--c shows
the cumulative distribution of these fractions over the reported top 100.
For PDZ trans-1, the concentration near zero means that some reported
leaders usually move down under the fitted analysis. The diagonal is a
descriptive guide, not a calibration target: exact rank recovery would
place every candidate at 0.5. Appendix~\ref{app:computation} gives the
formula and finite-sample conventions.

\paragraph{Direct rank concordance.}
Agreement varies across landscapes (Fig.~\ref{fig:reported}d--f). Under
Exponential--Exponential parameter resampling, the Spearman correlation
between reported and median sampled ranks in the reported top 100 is
0.943 for PDZ cis but 0.310 for PDZ trans-1. In the TrpB library-B
Exponential--Exponential example, the correlations are 0.050--0.056 across
the three schemes, compared with 0.503--0.584 for NucB generation 3 under
Exponential. Expanding the reported cohort to 1,000 increases these ranges
to 0.252--0.262 and 0.921--0.926, respectively
(Fig.~\ref{fig:reported1000}). Agreement over a larger leading cohort can
therefore hide disagreement within its top 100.

\begin{figure}[p]
\centering
\includegraphics[width=\linewidth]{reported_order.pdf}
\caption{Agreement with the reported top-100 ordering
(Section~\ref{sec:reported}). Panels (a)--(c) show empirical cumulative
distributions of reported-rank sample percentiles for PDZ trans-1 and TrpB
libB under Exponential--Exponential specifications and NucB G3
under the Exponential specification. Blue, red, and green denote parameter,
optimized-fit count, and posterior count resampling. The dashed diagonal
is descriptive, not a calibration target. Panels (d)--(f) show within-cohort
Spearman correlations for all analyzed landscapes. Small faint points denote
model runs; larger circles, squares, and triangles denote the median across
models within each landscape and scheme. The collections contain 49, 60,
and 49 runs, respectively.}
\label{fig:reported}
\end{figure}

\section{What is the impact of including noise?}
\label{sec:noise}
\subsection{Rank uncertainty within leading batches}
\label{sec:ranks}

Rank intervals for the reference top 100 often extend beyond the selected
batch. For TrpB library B under Exponential--Exponential parameter resampling,
the median central 90\% rank width is 281 positions, and 84 of the 100
candidates (84\%) have an upper endpoint beyond rank 100. For NucB
generation 3 under Exponential, the median width is 38 positions and
18 candidates (18\%) have an upper endpoint beyond 100.

Figure~\ref{fig:ranks}a--c displays intervals for the first 20 candidates
in each example's median-score reference. The blue points mark median
sampled ranks and the black ticks mark reference ranks; these can differ
because taking a median and ranking do not commute. Intervals can also
extend much farther toward worse ranks than toward better ones.

The broad rank intervals in TrpB library B persist across the noise
constructions examined (Fig.~\ref{fig:ranks}d--f): the median width among
the reference top 100 ranges from 281 to 860 positions, and 84--92 candidates
(84--92\%) have an upper endpoint above 100. These widths summarize
individual rank distributions, not an interval for an average rank.
Whether the rank movement changes batch membership depends on where it
occurs relative to the selection cutoff.

\begin{figure}[p]
\centering
\includegraphics[width=\linewidth]{rank_uncertainty.pdf}
\caption{Rank uncertainty within leading batches (Section~\ref{sec:ranks}). Panels
(a)--(c) show the first 20 candidates in the median-score ordering for the
PDZ trans-1 and TrpB libB Exponential--Exponential specifications and
the NucB G3 Exponential specification. Horizontal intervals
span the 5th--95th sampled-rank percentiles under parameter resampling; blue
points mark median sampled ranks and black ticks mark median-score reference ranks.
Horizontal scales differ across examples. Panels (d)--(f) summarize the median
interval width among each run's reference top 100. Faint points denote runs;
larger symbols denote medians across models within a landscape. Blue circles,
red squares, and green triangles denote the three resampling schemes. The
logarithmic horizontal axes report rank positions; the dashed line marks a
width of 100. These widths describe sample distributions, not validated
coverage against a known biological ordering.}
\label{fig:ranks}
\end{figure}

\subsection{Turnover in selected batches}
\label{sec:decisions}

For the TrpB library-B specification above, two resampled
top-100 batches share only 26--51 candidates (26--51\%) on average across
the three schemes (Fig.~\ref{fig:selection}e). Thus about 49--74 members
of either batch are replaced in the other, despite a fixed model and
selection budget within each analysis.

NucB generation 3 has much higher batch agreement: two sampled top-100
batches share 87--93 candidates (87--93\%) on average
(Fig.~\ref{fig:selection}f). Its uncertain internal ranks therefore coexist
with a comparatively stable candidate group. Table~\ref{tab:reference-comparisons}
gives the overlaps for each scheme. These quantities describe
the fitted resampling distributions, not observed agreement between new
physical experiments.

\paragraph{Candidates near the cutoff.}
Figure~\ref{fig:selection}a--c shows reference ranks
91--100 and the ten outside candidates with the largest inclusion probabilities
under parameter resampling. In TrpB library B, several outside candidates are
selected about as often as candidates just inside the reference top 100.
The cutoff therefore separates candidates with similar selection frequencies.
NucB generation 3 shows a clearer difference between the two groups.

\paragraph{Internal stability and agreement with a reference.}
Agreement between sampled batches need not imply agreement with the
reported batch. Under parameter resampling for NucB generation 3,
a sampled top-100 set shares 95.36 candidates on average with its median-score
reference but only 68.14 with the reported set
(Table~\ref{tab:reference-comparisons}). For TrpB library B, the corresponding
values are 63.44 and 29.81, while two sampled sets share 50.96 candidates.
NucB mainly differs from the reported list; TrpB also has substantial
variation within the fitted analysis. Overlap with the reported answer
alone would not distinguish these cases.

\begin{figure}[p]
\centering
\includegraphics[width=\linewidth]{selection_stability.pdf}
\caption{Turnover in selected batches (Section~\ref{sec:decisions}).
Panels (a)--(c) show top-100 inclusion probabilities under parameter resampling
for median-score reference ranks 91--100 (blue) and the ten highest-probability
outside candidates (orange). Outside candidates are ordered by decreasing
inclusion probability; their row numbers are not their sampled ranks. Panels
(d)--(f) show mean shared fractions between distinct sampled top-$k$ sets for
the three resampling schemes, with batch sizes $k$ from 10 to 1,000.
Each curve is one run, not an average over landscapes. PDZ trans-1 and TrpB libB
use Exponential--Exponential models; NucB G3 uses Exponential.
The example parameter and posterior-count runs have
16,000 samples and the optimized-fit count runs have 1,000. The vertical
guide marks $k=100$. A shared fraction of 0.5 means that two size-$k$ sets
share $k/2$ variants on average, not that the whole sets match half the time.}
\label{fig:selection}
\end{figure}

Batch size changes the conclusion, but increasing the budget does not always
increase the shared fraction.
After taking the median across model specifications within each landscape,
$\widehat C_{1000}>\widehat C_{100}$ in 13 of 18 landscapes (72\%) for
parameter resampling, 14 of 19 (74\%) for optimized-fit count resampling,
and 13 of 18 (72\%) for posterior count resampling. The same five TrpB
libraries, D through H, decrease under all three schemes
(Table~\ref{tab:landscape-set-sizes}). Adjusting overlap for the uniform-subset
reference using $(\widehat C_k-k/N)/(1-k/N)$ leaves the counts and the
five exceptions unchanged. This compares the endpoints 100 and 1,000;
it does not establish monotonicity between them.

\section{Why do we see these patterns?}
\label{sec:mechanisms}

Lower input counts are associated with wider score intervals in all 158
decision-cohort analyses. But that association alone does not say which
candidates can replace members of a selected batch. Figure~\ref{fig:explanation}
follows the counts, score uncertainty, and selection frequencies of every
reference member and challenger in three examples. Their different patterns
show why the candidate-level data matter alongside the overall correlations.

\paragraph{Count support and score precision.}
For each variant, we measure score width by
$W_i^{\fit}:=Q_{0.95}(\fit_i)-Q_{0.05}(\fit_i)$. To compare within-run
score geometry across different numerical scales, let $S_{\mathrm{run}}$
be the interquartile range (IQR) of the median scores among the reference top 1,000.
We use $U_i:=W_i^{\fit}/S_{\mathrm{run}}$ as normalized score uncertainty.
The denominator puts score widths on the scale of differences between
leading candidates. It is not an uncertainty estimate for one candidate.
Write $n_i^{\mathrm{in}}$ and $n_i^{\mathrm{out}}$ for total observed input
and output reads over the retained observations, with shared input counts
included once. Appendix~\ref{app:count-features} defines the sums for each
assay, including timepoints, gates, and resequencing.

Within each decision cohort, the Spearman correlation between input counts
and $U_i$ is negative in all 158 runs; the output-count correlation is
negative in 157. Figure~\ref{fig:counts} shows these correlations separately
by landscape and scheme. Multiple runs use the same experimental counts,
so the 158 results are not independent replications. Output counts also
depend on selection, whereas input counts are measured before it.
Appendix~\ref{app:per-analysis-why} shows the point-level relationships for
every run, without binning or fitting trend lines.

\paragraph{Locating plausible replacements.}
Let $\widetilde\fit_{(100)}$ be the score at the reference rank-100 cutoff.
The normalized margin
$X_i:=(\widetilde\fit_i-\widetilde\fit_{(100)})/S_{\mathrm{run}}$
locates each candidate relative to that cutoff.
Positive margins identify reference top-100 members, subject to exact ties;
negative margins identify candidates outside the reference set. Each row
of Fig.~\ref{fig:explanation} shows score width against input counts, output
counts, and this margin, followed by inclusion frequency against the same
margin. Filled blue points are reference members and open orange points
are challengers. The same candidates appear in all four panels of a row.

\paragraph{PDZ trans-1.}
Score width declines as either count increases in this example
(Fig.~\ref{fig:explanation}a,b), although the displayed widths span a
relatively narrow range on the common scale. Reference members and
challengers have overlapping widths (panel c). Selection frequency changes
gradually across the cutoff (panel d), and 319 outside candidates have
inclusion frequency at least 0.01. The proposed top 100 therefore has many
plausible replacements, even with a consistent count--width relationship.

\paragraph{TrpB library B.}
The count--width relationship spans a much wider range here
(Fig.~\ref{fig:explanation}e,f). Many challengers with median scores below
the cutoff have wider score intervals than reference members (panel g).
Some lie several run-IQR units below the cutoff but still enter the
top 100 in at least 1\% of samples (panel h). There are 416 challengers
in total. Selecting only from the median-score ordering would leave these
candidates out of the proposed batch, despite their nonzero inclusion
frequencies under resampling.

\paragraph{NucB generation 3.}
The pattern differs in NucB (Fig.~\ref{fig:explanation}i--l). Many reference
members have wider score intervals than the challengers, reversing the
margin--width pattern seen in TrpB (panel k). Yet 77 of the 100 reference
members (77\%) have inclusion frequency at least 0.99, and only 28 outside
candidates meet the challenger threshold (panel l). A broad individual
score interval does not by itself imply uncertain batch membership.
Inclusion depends on the relative scores in each joint sample, not on
the width of one candidate's interval alone.

These examples connect the batch overlaps in Fig.~\ref{fig:selection}
with the candidates that produce them. Count support, score margins, and
inclusion frequencies describe different parts of that result; none can
replace the others. The assays and fitted models differ between rows, so
the comparison does not isolate an effect of counts or score spacing.

We also fit descriptive regressions of rank width on score width, local
crowding, their interaction, and reference rank. Local crowding is the
negative logarithm of normalized score spacing, so larger values mean
closer competitors. The coefficients vary across analyses, and we do not
find a consistently positive interaction. Appendix~\ref{app:associations}
reports the specifications and coefficient counts. They summarize these
fits; they do not estimate how extra sequencing would change the ranks.

\begin{figure}[p]
\centering
\includegraphics[width=\linewidth]{explaining_resolution.pdf}
\caption{Count support and plausible batch replacements
(Section~\ref{sec:mechanisms}). Rows show PDZ trans-1 and TrpB libB
under Exponential--Exponential, and NucB G3 under Exponential,
each with 16,000 parameter samples. Every reference top-100 member (blue,
filled) and outside challenger with inclusion frequency at least 0.01
(orange, open) is retained: 419, 516, and 128 candidates, respectively.
From left to right, columns show central 90\% score width against input
count, output count, and median-score margin from the top-100 cutoff;
then top-100 inclusion frequency against that margin. Widths and margins
are divided by the run's top-1,000 median-score interquartile range (IQR).
Count axes show $\log_{10}(n+1)$ and width axes are logarithmic. Count and
width limits are shared across rows; margin limits are shared within a row
but differ between rows. Dotted lines mark zero margin. No points are
binned or subsampled, and no regression lines are fitted.
Appendix~\ref{app:figure-explanation} gives the example and display conventions.}
\label{fig:explanation}
\end{figure}

\paragraph{What normalization does and does not remove.}
Dividing by $S_{\mathrm{run}}$ does not change within-run Spearman
correlations. Nonlinear monotone score transformations can, however, change
widths and spacing while preserving every sampled ranking. We therefore
interpret these score summaries on their stated scales and use sampled
ranks and selections to measure the ranking consequences.

\paragraph{Dependence on modeling assumptions.}
\label{sec:sensitivity}

For the example specifications, TrpB has lower batch agreement than NucB
under all three sampling recipes. Other comparisons are much more sensitive
to model choice. In PDZ trans-1, parameter-resampling top-100 overlap is
59.80\% under Exponential--Exponential and 0.27\% under Gamma--Exponential.
We cannot interpret the larger uncertainty as a property of the experiment
without checking the corresponding model. Appendix~\ref{app:per-analysis-why}
shows all runs separately, and Appendix~\ref{app:fit-checks} reports the
available diagnostics and missing predictive checks.

\section{Discussion}
\label{sec:discussion}

We found that fitness calculation and measurement uncertainty can change
the composition of leading batches in deep mutational scans.

\paragraph{Limitations.}
The estimated uncertainty is conditional on the assay models and retained
fits. A model that does not reproduce the data can give misleading rank
intervals and inclusion frequencies. Predictive checks remain incomplete;
six posterior fits have tree-depth saturation, and nine lack stored
native-score convergence summaries (Appendix~\ref{app:fit-checks}). The
different sampling recipes also share modeling assumptions, so agreement
between them is not independent validation of an experiment's uncertainty.
The studied landscapes are not a representative sample of all scans.

\paragraph{Protein engineering.}
An engineer choosing a batch can use inclusion frequencies to distinguish
candidates that are repeatedly selected from those with plausible replacements.
The TrpB example shows that these replacements can involve much of the
batch. But turnover alone is not a loss in engineering utility: the
alternative variants may have similar activity. Comparing their activity
or the expected outcome of a specified follow-up experiment is a separate
question. The median-score list is a reference for this analysis, not an
optimal selection rule for every experimental objective.

\paragraph{Machine-learning benchmarks.}
Fitness scores serve both as training labels and as references for
evaluating predictions. Here we quantify uncertainty in the leading
experimental ranks and batches, not its effect on any predictor. For
landscapes used in FLIP2, a direct next step is to hold predictions fixed
and evaluate them across sampled score sets within each task's candidate
subset. This would measure how label uncertainty affects rank correlation,
retrieval, and comparisons between predictors.

\paragraph{Further checks.}
Two kinds of evidence would strengthen the interpretation. First, predictive
checks and suitable held-out measurements could test whether the fitted
models reproduce the relevant experimental variation. Second, controlled
read thinning or leave-one-replicate-out fits could test whether changing
count support changes batch stability. Both require analyses of altered
or additional observations; subsampling existing model draws only assesses
Monte Carlo variability.

The current results concern candidate identity at a specified batch size.
For the TrpB example, a top-100 list leaves many plausible replacements;
for NucB, much more of that list remains selected. Reporting those
selection frequencies alongside fitness scores tells a reader what an
exact ordering leaves uncertain.

\section*{Data and code availability}
PDZ processed counts and scores are available through GEO accession GSE265816.
TrpB data and analysis code are deposited at
\url{https://doi.org/10.22002/h5rah-5z170}. NucB data and source calculations
are available from \url{https://github.com/google-deepmind/nuclease_design}.
The rank-analysis code consumes prepared sampled-score matrices and aligned
count features; fitting the original models is not required to recreate its
tables and figures. The accompanying repository includes a small pipeline
example, which is distinct from the real-data analyses reported here.
\note{TODO: Add versioned fitting and analysis repositories and a public archive
for the prepared samples, features, and fit diagnostics.}

\section*{Author contributions}
\note{TODO: Confirm contributions with all authors.}

\section*{Competing interests}
\note{TODO: Obtain declarations from all authors.}

\section*{Acknowledgments}
\note{TODO: Confirm funding and acknowledgments.}

\clearpage
\input{paper/old_paper_ml_structure/references.tex}

\clearpage
\appendix
\renewcommand{\thefigure}{S\arabic{figure}}
\renewcommand{\theHfigure}{supp.\arabic{figure}}
\setcounter{figure}{0}

\input{paper/old_paper_ml_structure/supplementary_details.tex}

\section{Count-based fitness calculations}
\label{app:scores}

\paragraph{PDZ.}
Let $n_{ir}^{\mathrm{in}}$ and $n_{ir}^{\mathrm{out}}$ denote counts for
variant $i$ in replicate $r$, and let $w$ index the wild-type reference. Let
$\mathcal R_i$ be the set of replicates for which all four counts for $i$
and $w$ are positive. For $r\in\mathcal R_i$, define the wild-type-relative
log enrichment $f_{ir}$ and its working variance $v_{ir}$ by
\begin{align}
f_{ir}&:=\log\frac{n_{ir}^{\mathrm{out}}}{n_{ir}^{\mathrm{in}}}
-\log\frac{n_{wr}^{\mathrm{out}}}{n_{wr}^{\mathrm{in}}},\\
v_{ir}&:=\frac{1}{n_{ir}^{\mathrm{in}}}+\frac{1}{n_{ir}^{\mathrm{out}}}
+\frac{1}{n_{wr}^{\mathrm{in}}}+\frac{1}{n_{wr}^{\mathrm{out}}},\qquad
g_i(\data):=\frac{\sum_{r\in\mathcal R_i} f_{ir}/v_{ir}}
{\sum_{r\in\mathcal R_i}1/v_{ir}}.
\end{align}
Only valid replicates enter the sums. We define the score as undefined when
$\mathcal R_i$ is empty, so undefined logarithms never enter the weighted
average. This inverse-Poisson-variance
calculation is a DiMSum-style approximation. The exact source DiMSum revision
and configuration have not been recovered, so we do not describe it as an
exact reproduction of the published score \citep{faure2020,zarin2024}.

\paragraph{TrpB.}
Following the deposited calculations of \citet{johnston2024}, we form log
output/input ratios separately for each retained timepoint and replicate.
We subtract the same-observation mean for stop-codon variants and divide by
the difference between the same-observation maximum and that stop-variant
mean. The final score averages valid normalized observations. The three-site
libraries use one input count per variant; the four-site landscape uses the
mean of two input-replicate counts. Inputs must be at least 10 and required
outputs positive; if an observation has two output replicates, both must be
positive. Library-specific timepoints follow the deposited notebooks; the
four-site calculation excludes 12 hours and retains 16, 20, 24, 28, and
36 hours. Culture-density and time factors cancel from the within-observation
stop/maximum normalization. The source count calculations reproduce the
published finite scores and coverage at numerical tolerance $10^{-12}$.

\paragraph{NucB.}
For the highest selected fluorescence gate, we calculate
\[
g_i(\data):=\frac{n_i^{\mathrm{gate}}/D^{\mathrm{gate}}}
{n_i^{\mathrm{in}}/D^{\mathrm{in}}},
\]
where $D^{\mathrm{gate}}$ and $D^{\mathrm{in}}$ are the source normalization
totals, including controls when specified by the source. The selected gates
are generation-specific \citep{thomas2025}. A zero gate count gives zero
enrichment; a zero input count makes the score undefined. Source controls
contribute to the appropriate totals but not to ranking candidates. Generation
1 retains the source total correction for sequences omitted from the fitted
axis; this common multiplicative correction does not alter ranks. The observed
count calculations reproduce the source enrichment values at tolerance
$10^{-12}$.

\section{Ranking and computational conventions}
\label{app:computation}

\paragraph{Finite samples and reference support.}
We store sampled scores as single-precision arrays and retain samples with at
least 1,000 finite candidate scores. Table~\ref{tab:inventory} gives the
resulting sample-count ranges. Rank intervals and reported-rank percentiles
use the finite observations for each candidate; undefined scores do not become
zero scores. Inclusion counts use membership in the retained size-$k$ sets,
with the retained sample count as denominator. Count-based median-score references
require at least two finite scores per candidate; the parameter collections
require at least one. The stored sample identities preserve alignment across
parameters and count draws. Source controls and other excluded candidates
remain outside the ranking universe.

\paragraph{Ranks and selections.}
Exact tied scores receive average ranks for interval summaries. Fixed-size
top-$k$ selection uses the deterministic selection order, with source block
order retained where needed to reproduce the prepared results. Because the
two conventions answer different questions at a tie, we calculate set
membership from stored selections rather than from average rank thresholds.
The median-score reference orders median scores and resolves its exact ties
deterministically. Rankings always use the eligible full competitor pool;
the reported cohorts and displayed summaries remain top-focused.

\paragraph{Exact shared fractions.}
Let $c_i(k):=\sum_s\mathbf{1}\{i\in T_k^{(s)}\}$ be the number of retained
sets containing variant $i$. Then the sample-to-sample overlap above is
exactly
\[
\widehat C_k=\frac{\sum_i c_i(k)[c_i(k)-1]}{S(S-1)k}.
\]
We use this identity instead of drawing random pairs of samples. It is an
empirical summary of distinct retained samples; dependence among Monte Carlo
draws is not removed by the identity. We likewise obtain reference overlap
by averaging inclusion frequencies over each fixed reference set.

\paragraph{Reported-rank percentiles.}
Let $r_i^{\mathrm{rep}}$ be candidate $i$'s reported rank and $\mathcal S_i$
the indices of its finite sampled ranks. The fraction plotted in
Fig.~\ref{fig:reported} is
\[
u_i:=\frac{1}{|\mathcal S_i|}\sum_{s\in\mathcal S_i}
\left[\mathbf1\{R_i^{(s)}<r_i^{\mathrm{rep}}\}
+\tfrac12\mathbf1\{R_i^{(s)}=r_i^{\mathrm{rep}}\}\right].
\]
We plot its empirical cumulative distribution over the reported top 100
or top 1,000, omitting candidates with no finite sampled rank. Exact rank
recovery gives $u_i=0.5$ for every candidate, so a uniform diagonal is not
a calibration target. These percentiles and the reported-versus-median-rank
correlations describe agreement with a fixed reference, not errors against
a known biological ranking.

\paragraph{Implementation checks.}
The repository tests compare streamed rankings with dense calculations,
exercise ties and candidate exclusions, and check table and figure outputs.
These are checks of the analysis implementation, not a demonstration of
frequentist coverage or biological model adequacy. A controlled synthetic
coverage study is not part of the present draft.

\section{Descriptive association analyses}
\label{app:associations}

The unaggregated relationships behind this section are shown for every
landscape, model, and resampling scheme in
Appendix~\ref{app:per-analysis-why}. Its scatterplots contain no fitted
regression lines. The regression coefficients summarized here are separate
descriptive calculations; the point displays let the reader assess patterns
that a single coefficient or an across-landscape median can obscure.

Recall that $U_i$ is normalized score width and $\widetilde r_i$ is the
median-score reference rank. Local score spacing is the average change in
median score per rank across up to five neighbors on each side of a reference
rank, divided by $S_{\mathrm{run}}$. Endpoints use the available neighbors.
We define crowding as minus the base-10 logarithm of this spacing, using a
small numerical floor for zero spacings. Larger values indicate closer
competitors on this scale.

\paragraph{Figure 4 examples and plotted data.}
\label{app:figure-explanation}
Figure~\ref{fig:explanation} uses the parameter-resampling versions of the
same three examples as Figs.~\ref{fig:reported}--\ref{fig:selection}:
PDZ trans-1 and TrpB library B under Exponential--Exponential, and NucB
generation 3 under Exponential. No Gamma specification enters this figure.
The examples illustrate different candidate-level patterns; they are not
a representative sample of model specifications. The complete per-analysis
displays, including Gamma fits, remain in Appendix~\ref{app:per-analysis-why}.

Each row includes every member of the median-score reference top 100 and
every outside candidate with $\hat q_i(100)\geq0.01$. The three rows contain
1,063 candidates in total. Every candidate appears in all four panels of
its row, without binning, point subsampling, or pooling across analyses.
Each point summarizes one candidate's score samples; it is not an
individual sample. The stored input counts, output counts, score medians,
interval endpoints, and inclusion frequencies are unchanged.

Counts are transformed as $\log_{10}(n+1)$, retaining zero counts at zero.
Normalized widths use a logarithmic axis. All widths in these three examples
are positive, so no plot-only floor affects the displayed values. Count
and width limits cover all points and are common to the three rows.
The two margin panels use the same limits within a row, but the limits
differ between rows; apparent transition steepness should therefore not
be compared without reading the axes. Inclusion frequencies use a common
zero-to-one scale.

The exported variant table retains every plotted row and its coordinates,
and the run index records example identities, cohort sizes, sample counts,
and mean shared fractions. Inclusion frequencies and shared fractions come
from the joint sampled rankings against all finite eligible competitors,
not from a candidate's marginal interval or a ranking restricted to the
displayed cohort. Because challengers are selected using their inclusion
frequencies, relationships in this cohort do not describe all variants
in the landscape.

\paragraph{Observed count features.}
\label{app:count-features}
The features are fixed from the observed data and shared across resampling
schemes for a landscape. PDZ input and output support sum the three paired
replicates separately. TrpB input support uses the single input count once
for three-site libraries and the sum of the two input counts for four-site;
output support sums every retained output replicate and model timepoint.
NucB input support sums the arrays \texttt{ic1}, \texttt{ic2}, and
\texttt{ic3} where present, over all their observation axes. Output support
sums the low- and high-gate arrays for G1, the three gate arrays for G2,
and all four output rows for G3 and G4. Resequenced libraries contribute
their recorded reads to these descriptive totals. None of these sums turns
sequencing repeats, different timepoints, or different gates into independent
biological replicates. Log input support means $\log(1+n_i^{\mathrm{in}})$;
the count scatterplots use $\log_{10}(1+n)$ for display.

\paragraph{Observation discordance.}
For PDZ paired replicates, let $e_{ir}$ be the difference of smoothed log
output and input frequencies for variant $i$ and replicate $r$. Specifically,
\[
e_{ir}:=\log\frac{n_{ir}^{\mathrm{out}}+\delta}
{D_r^{\mathrm{out}}+M\delta}
-\log\frac{n_{ir}^{\mathrm{in}}+\delta}{D_r^{\mathrm{in}}+M\delta},
\qquad
v_{ir}^{\mathrm{obs}}:=\frac{1}{n_{ir}^{\mathrm{in}}+\delta}
+\frac{1}{n_{ir}^{\mathrm{out}}+\delta},
\]
where $\delta=0.5$ and $D_r$ are totals on the feature-construction count
axis of length $M$. Set $w_{ir}:=1/v_{ir}^{\mathrm{obs}}$ and
$\overline e_i:=\sum_rw_{ir}e_{ir}/\sum_rw_{ir}$. The reduced
heterogeneity is
\[
H_i:=\frac{1}{R-1}\sum_{r=1}^R w_{ir}(e_{ir}-\overline e_i)^2,
\]
where $R=3$ for PDZ. It measures replicate scatter relative to a working
Poisson log-count variance. The smoothing is part of this diagnostic feature;
it does not change the no-pseudocount fitness reconstructions in
Appendix~\ref{app:scores}.

For TrpB observations with two output replicates, we instead use smoothed log
output frequencies and variance $1/(n_{i\ell r}^{\mathrm{out}}+\delta)$,
remove a separate weighted mean within each timepoint $\ell$, and sum the
weighted squared residuals over timepoints. With $J$ timepoints and $R=2$
replicates, the denominator is $J(R-1)$. This removes between-timepoint
changes before measuring within-timepoint disagreement. TrpB A--C have no
paired output contrast; NucB has no verified comparable feature in the
current analysis. Their discordance value is missing, not zero. The regression
includes this predictor only when more than three decision-cohort values
are available, then uses complete cases.

Within the reference top-100-plus-challengers cohort, the score-precision
regression uses $\log(1+U_i)$ as response and log input support and
$\log(1+\widetilde r_i)$ as predictors, adding the count-aware observation
discordance feature where available. The rank regression uses
$\log[1+Q_{0.95}(R_i)-Q_{0.05}(R_i)]$ as response and $U_i$, crowding,
their product, and $\log(1+\widetilde r_i)$ as predictors. We fit linear
regressions after standardizing the response and each predictor on complete
cases. Coefficients are summarized within landscapes by their median across
models. Their signs describe these fitted associations; the interaction
means that a main-effect coefficient is not a universal effect at every
value of the other predictor.

The score-uncertainty coefficient is positive in 17 of 18, 18 of 19,
and 17 of 18 landscape summaries for parameter, optimized-fit count, and
posterior count resampling, respectively. The crowding coefficient is
positive in 15 of 18, 17 of 19, and 14 of 18. The interaction is positive
in 3 of 18, 7 of 19, and 6 of 18. These are counts of fitted coefficient
signs, not uniformly positive conditional slopes. The predictors and rank
widths summarize the same samples, so the regressions are descriptive
rather than independent tests of an effect of additional sequencing.

\section{Supporting figure views}

\input{paper/old_paper_ml_structure/figures/landscape_set_sizes.tex}

\begin{figure}[p]
\centering
\includegraphics[width=\linewidth]{reported_order_top1000.pdf}
\caption{Agreement within the reported top-1,000 cohort. Panels (a)--(c)
show reported-rank percentile distributions for PDZ trans-1 and TrpB libB
under Exponential--Exponential and NucB G3 under Exponential.
Blue, red, and green denote parameter, optimized-fit count, and posterior
count resampling. Panels (d)--(f) show reported-versus-median-sampled-rank
Spearman correlations across all analyzed landscapes; faint points are model
runs and larger symbols are medians across models within each landscape.
The displayed ranks are still ranks against all finite eligible
competitors, not ranks recalculated only within the cohort.}
\label{fig:reported1000}
\end{figure}

\begin{figure}[p]
\centering
\includegraphics[width=\linewidth]{count_support.pdf}
\caption{Count support and score uncertainty across all analyzed landscapes.
Panels show within-run Spearman correlations of input counts (left) and output
counts (right) with normalized central 90\% score width. Each run uses its
reference top 100 plus challengers with inclusion probability at least 0.01.
Faint points are model runs; larger symbols are medians across models within
landscape and scheme. Blue circles, red squares, and green triangles represent
parameter, optimized-fit count, and posterior count resampling. These
are associations in selected cohorts, not independent experimental effects.}
\label{fig:counts}
\end{figure}

\clearpage
\section{Per-analysis count support and selection uncertainty}
\label{app:per-analysis-why}

We show each of the 158 analysis runs separately, grouped by protein system
and resampling scheme. Table~\ref{tab:why-index} locates the figures; the
row headings give the landscape, model specification, retained sample count,
decision-cohort size, and mean shared top-100 fraction. Each row contains the
median-score reference top 100 and all outside candidates with top-100
inclusion frequency at least 0.01. These are the same top-focused cohorts
used in Section~\ref{sec:mechanisms}, not samples of all eligible variants.

From left to right, the panels show normalized score width against input
support, output support, and the score margin from the top-100 cutoff;
the last panel shows top-100 inclusion against that same margin. Counts and
score widths are defined in Appendix~\ref{app:associations}. The score-width
and margin denominator is each run's interquartile range (IQR) of the
top-1,000 median scores. Count axes display $\log_{10}(1+n)$, and score
widths use logarithmic axes. Count limits and score-width limits are common
within each protein-system and scheme collection, including across pages.
The two margin panels share a scale within a row; their scales can differ
between runs because their score ranges differ.

All points are retained. Exact-zero score widths are displayed at half the
smallest positive normalized width in the corresponding protein-system and
scheme collection; the numerical tables retain zero. The accompanying figure
index records those plot-only floors and the number of affected points.
Blue filled points denote reference top-100 members and orange open points
denote challengers. The dotted vertical lines mark a score margin of zero.
No regression lines or smoothing curves are fitted. These figures display
within-analysis associations, not causal effects of changing read depth and
not model-adequacy checks against new experimental observations.

\input{paper/old_paper_ml_structure/figures/summary_why.tex}

\end{document}